\textbf{Neural execution and benchmarks.} Veli\v{c}kovi\'c et al.\ \cite{velikovi2019neural} trained message-passing networks to imitate breadth-first search, Bellman-Ford and Prim step by step and argued for max aggregation; our max-aggregation, step-supervised system is a reimplementation of that recipe in spirit (our architecture and losses are our own simplified version, not their code). The CLRS benchmark \cite{velikovi2022clrs} standardises 30 algorithms with hints and a 16-to-64 node size split, the triplet-GMPNN processor \cite{ibarz2022generalist} and pointer graph networks \cite{velikovi2020pointer} are other processors that we do not reimplement, and SALSA-CLRS \cite{minder2023salsa} supplies sparse large-graph versions. We use a smaller 8-to-64 split on one algorithm.

\textbf{Why aggregation should matter.} Xu et al.\ \cite{xu2019what} show that networks whose computation graph is close to the algorithm's dynamic program need fewer samples, and Xu et al.\ \cite{xu2020how} show that ReLU MLPs extrapolate linearly outside the training range and that GNNs with min/max aggregation can extrapolate on shortest-path-like tasks when the architecture matches. Dudzik and Veli\v{c}kovi\'c \cite{dudzik2022graph} relate GNNs to dynamic programming and discuss why max aggregation fits it. The message-passing framework itself is from Gilmer et al.\ \cite{gilmer2017neural}, in the broader graph-network view of Battaglia et al.\ \cite{battaglia2018relational}.

\textbf{Size generalisation.} Yehudai et al.\ \cite{yehudai2020from} show that GNNs can fail to generalise to larger graphs because local neighbourhood structures (such as degree) differ between small and large graphs, which motivates our two graph families. Bevilacqua et al.\ \cite{bevilacqua2021size} study size-invariant representations for graph classification, and Bevilacqua et al.\ later propose a causal self-supervised regulariser that improves out-of-distribution size generalisation on CLRS tasks \cite{bevilacqua2023neural}. We evaluate none of these remedies.
